\vspace*{-3ex}
\subsection{Adaptive Similarity Measure}
\vspace*{-1ex}
\label{sec-dpf}

Using some results of our extensive image-data mining \cite{DPF-ICIP02},
we formulate the {\em dynamic partial} distance function (DPF) as follows.
Let $\delta_i = |x_i-y_i|$ denote the difference 
between objects $x$ and $y$ in feature $i$, for $i=1,\cdots, p$.
We first define set $\Delta_m$ as
\[
\Delta_m = \{ The~~smallest~~m~\delta's~~of~~(\delta_1, ..., \delta_p)\}.
\]
The first approximation of a desired perceptual function looks like
\begin{equation}
d_{DPF}(X, Y) = (\sum\limits_{\delta_i \in \Delta_m}{\delta_i}^{r})^{\frac{1}{r}
}.
\label{eq-dpf}
\end{equation}
The above formulation shows that
DPF activates different features for
different object pairs. The activated features are those
with minimum differences---those which provide coherence
between the objects.  If coherence can be maintained
(because a sufficient number of features are similar), then
the objects paired are perceived as similar.
Our empirical studies in image retrieval \cite{DPF-ICIP02}
and in shot-boundary detection \cite{WC02} have shown that
DPF performs significantly better than Minkowski-like
functions.  

The effectiveness of DPF can be explained in
similarity theories in cognitive science.
Some important work in the cognitive psychology community
has provided ample evidence that
similarity is both {\em dynamic} and {\em non-metric}
\cite{MM85, Tev77}.  Most recently, Goldstone \cite{Gold93}
shows that similarity perception
is the {\em process} that determines the respects (i.e.,
features, or attributes) for measuring similarity.
More precisely, a similarity function for measuring
a pair of objects is formulated only after the objects
are compared, not before the comparison is made; and the
respects for the comparison are {\em dynamically} activated in
this formulation process.  Let us use a simple example to explain.
Suppose we are asked to name two places that are similar
to England.  Among several possibilities, Scotland and
New England could be two reasonable answers.  However,
the respects in which England is similar to Scotland differ from those
in which England is similar to New England.
If we use the shared attributes of
England and Scotland to compare England and New England,
the latter pair might not seem similar, and vice versa.
Thus, a distance function using a fixed set of respects
(e.g., the traditional weighted Minkowski function) cannot
capture objects that are similar in different sets of respects.

\subsubsection*{\underline{Task: Perceptual Distance Functions}}

As we have discussed in Section~\ref{sec-dpf}, DPF dynamically
selects $m$ out of $p$ total features for measuring
similarity between two objects.
Even though DPF performs markedly better than traditional
distance functions, there is still much room for improvement.  

When the similarity of an object-pair is measured, DPF 
chooses $m$ feature-channels that have the smallest distance
to produce an aggregated distance.   Ideally, the optimal $m$ value should
be determined in a pairwise fashion.  But the current version 
of DPF uses an ``on average optimal'' $m$ learned from
a set of training pairs.  This one-$m$-fits-all approach
clearly cannot work optimally for all object-pairs.

We plan to investigate two methods to tune $m$ adaptively for each 
object-pair: {\em thresholding} and {\em sampling}.
\begin{nlist}
\item {\bf 1. \em Thresholding}.
 Intuitively, we know that we need to find $m$ similar feature-channels 
 between two objects.  The thresholding method finds for
 each object-pair $m$ channels that have a distance 
 below a given threshold.  In other words, given a threshold,
 the $m$ value can be determined pairwise.
 A potential shortcoming of this method is that
 we still need to find an ``on average
 optimal'' threshold for all object pairs.  
 
\item {\bf 2. \em Sampling}.
 Another method we plan to study is sampling,
 which picks a number of candidate $m$'s
 in a range, say, within three standard deviations of the
 ``on average optimal'' $m$. Using $w$ sampling points, DPF
 can compute $w$ different distance measures for each object pair.
 To return top-$k$ results, the sampling method can return 
 the objects that enjoy the highest ``votes'' among these $w$ 
 measures.

 The key idea behind the sampling method is that 
 among the $w$ sampling points, one of them is 
 likely to be near the optimal $m$.  We plan to
 conduct extensive experiments to find the best
 sampling range, sampling frequency, and voting method.
\end{nlist}

